In the present section we employ the circle decomposition theorem to study the decoration sequence, most importantly for $\L$-spectra. We will to this end extend the construction to arbitrary functors and analyse the resulting sequence for Grothendieck-Witt spectra using hermitian cofinality. This will result in one proof of Williams' conjecture from the introduction. The second half of the section then analyses the decoration construction more broadly and show that it, alongside the Lattice-Tate-construction already employed in the construction of $\KGW$ in Proposition \reffour{proposition:localising-are-stable}, in the limit gives rise to Karoubi-localising approximations for arbitrary Verdier-localising functors (and even a few more), without a detailed analysis of the entire decoration sequence, in particular resulting in a second proof of Williams' conjecture.



\subsection{Ultimate decorations and Williams' conjecture}
\label{subsection:universal-dec}%

The purpose of this subsection is to introduce the decoration construction and then use the circle decomposition formula established in \S\reffour{subsection:shaneson} together with hermitian cofinality in order to give a formula for $\KGW$ and $\KL$. In the case of $\KL$, this reproduces what is classically known as $\L$-theory with ultimate decoration, thus showing that the latter coincides with Karoubi $\L$-theory, in particular establishing a Williams' conjecture.  



 

\begin{definition}
\label{definition:decoration}%
Let $\F \colon \Catp \to \A$ be a functor valued in a stable $\infty$-category $\A$. We define $\T\F\colon \Catp \to \A$ by the formula
\[
\T\F(\C,\QF) = \Sig\TotFib\left[
\begin{tikzcd}
\F(\Hyp(\C)) \ar[r]\ar[d] & \F(\Hyp(\C_{\B\NN})) \ar[d] \\
\F(\C,\QF) \ar[r] & \F((\C,\QF)_{\B\ZZ})
\end{tikzcd}\right] \ ,
\]
where $\TotFib$ denotes the total fibre of the given square. $\T\F$ is therefore the cofibre of the induced map between the horizontal fibres. There is a natural map $\F \Rightarrow \T\F$ given as a composite
\[
\F(\C,\QF) = \fib[\F(\C,\QF) \to \B\ZZ \otimes \F(\C,\QF)] \Rightarrow \fib[\F(\C,\QF) \Rightarrow \F((\C,\QF)_{\B\ZZ})] \Rightarrow \T\F(\C,\QF).
\]
For $n \geq 0$ we then define $\F^{\cotower{-n}}\colon \Catp \to \A$ inductively as follows. We set $\F^{\cotower{0}} = \T\F$, and for $n \geq 1$ we set $\F^{\cotower{-n}} = \T\F^{\cotower{-n+1}}$. We then set
\[
\F^{\cotower{-\infty}} = \colim[\F \Rightarrow \F^{\cotower{0}} \Rightarrow \F^{\cotower{-1}} \Rightarrow ... ].
\]
We will refer to this as the \defi{decoration sequence} of $\F$, and to $\F^{\cotower{-\infty}}$ as \defi{ultimately decorated} approximation of $\F$. 
\end{definition}




Applied to $\F = \K$ this reproduces the Bass delooping of the $\K$-spectrum functor, so that $\K^{\cotower{-\infty}} = \KK$, see Lemma \reffour{proposition:connective-covers} below.


\begin{lemma}
\label{lemma:preserves-limits}%
Let $\A$ be a stable $\infty$-category. Then the following hold:
\begin{enumerate}
\item
\label{item:T-preserves-limits-and-colimits}%
The functor $\T \colon \Fun(\Catp,\A) \to \Fun(\Catp,\A)$ preserves additive functors as well as all limits and all colimits that exist in $\A$.
\item
\label{item:T-commutes-translation}%
There is an equivalence
\[
m_{(\D,\QFD)}^*\T(\F) \simeq \T(m_{(\D,\QFD)}^*\F)
\]
natural in $(\D,\QFD) \in \Catp$ and $\F\in \Fun(\Catp,\A)$, where $m_{(\D,\QFD)}^*$ is the endo-functor of $\Fun(\Catp,\A)$ given by $m_{(\D,\QFD)}^*\F \simeq \F((\D,\QFD) \otimes -)$. \item
\label{item:TFhyp-equivalence}%
There is a natural equivalence $(\T\F)^{\hyp} \simeq \T(\F^{\hyp})$ in $\Fun(\Catp,\A^{\B\Ct})$. 
\end{enumerate}
\end{lemma}
\begin{proof}
The first half of~\reffouritem{item:T-preserves-limits-and-colimits} follows directly from Proposition~\reftwo{lemma:Morita-exact-criterion}. The other statement follow from the following analysis: 
Using the projection formula from Corollary~\refone{corollary:hyp-projection} and the compatibility of tensors and the monoidal product from Remark~\refone{remark:comparing-2}, we may equivalently write the functor $\T\colon \Fun(\Catp,\A) \to \Fun(\Catp,\A)$ as
\[
\T\F(-) = \Sig\TotFib\left[
\begin{tikzcd}
\F((-) \otimes \Hyp(\Spaf)) \ar[r]\ar[d] & \F((-) \otimes \Hyp(\Spaf_{\B\NN})) \ar[d] \\
\F((-) \otimes (\Spaf,\QF^{\uni})) \ar[r] & \F((-) \otimes (\Spaf,\QF^{\uni})_{\B\ZZ})
\end{tikzcd}\right] \ .
\]
More precisely, $\T$ is the suspension of the total fibre of the square $\tau\colon \Del^1 \times \Del^1 \to \Fun(\Catp,\A)$ which is the image of the square 
\[
\begin{tikzcd}
\Hyp(\Spaf) \ar[r]\ar[d] & \Hyp(\Spaf_{\B\NN}) \ar[d] \\
(\Spaf,\QF^{\uni}) \ar[r] & (\Spaf,\QF^{\uni})_{\B\ZZ})
\end{tikzcd}
\]
under the map
\[
\Catp \to \Fun(\Fun(\Catp,\A),\Fun(\Catp,\A)) \quad\quad (\D,\QFD) \mapsto m_{(\D,\QFD)}^*
\]
Since each of the functors $m_{(\D,\QFD)}^*$ preserves limits and colimits and $\A$ is stable we conclude that $\T$ preserves limits and colimits, showing the second half of~\reffouritem{item:T-preserves-limits-and-colimits}. The equivalence $m_{(\Dtwo,\QFDtwo)} \circ m_{(\D,\QFD)} = m_{(\D,\QFD)}\circ m_{(\Dtwo,\QFDtwo)}$ determined by the symmetric monoidal structure of $\Catp$ then induces the equivalence of~\reffouritem{item:T-commutes-translation}. Finally, \reffouritem{item:TFhyp-equivalence} is essentially a special case of \reffouritem{item:T-commutes-translation} for $(\D,\QFD) = \Hyp(\Spaf)$, using again the projection formula from Corollary~\refone{corollary:hyp-projection}), only that the result is interpreted as taking values in $\A^{\B\Ct}=\Fun(\B\Ct,\A)$ via the $\Ct$-action on $\Hyp(\Spaf)$ as a Poincaré $\infty$-category. 
\end{proof}



\begin{proposition}
\label{proposition:T-of-localising}%
Let $\F\colon \Catp \to \A$ be a Verdier-localising functor valued in a stable $\infty$-category $\A$. Then $\T\F$ is naturally equivalent, as a functor under $\F$, to the functor $(\C,\QF) \mapsto \F(\C^\natural,\QF)$. In particular, if $\F$ is Karoubi-localising then the map $\F \Rightarrow \T\F$ is an equivalence.
\end{proposition}
\begin{proof}
Consider the commutative diagram
\[
\begin{tikzcd}
\F\Hyp(\C)\ar[r,equal]\ar[d]  & \F\Hyp(\C) \ar[r]\ar[d] & \F(\Hyp(\C_{\B\NN}))\ar[d] \\
\F(\C,\QF) \ar[r] & {\B\ZZ} \otimes^{\flat} \F(\C,\QF) \ar[r, "\ass^{\flat}"] & \F((\C,\QF)_{\B\ZZ}) \ .
\end{tikzcd}
\]
By Theorem~\reffour{theorem:circle-decomposition} and Remark~\reffour{remark:circle-decompose-2} the right square is exact, and hence the total cofibre of the external rectangle is naturally equivalent to the total cofibre of the left square, which in turn is given by $\cof[\F(\C,\QF) \to {\B\ZZ} \otimes^{\flat} \F(\C,\QF)] = \Sig\F(\C^\natural,\QF)$. %
We conclude that $\T\F(\C,\QF) = \F(\C^{\natural},\QF)$, as desired.
\end{proof}



Suppose that $\F\colon \Catp \to \A$ is a functor and $\F \Rightarrow \FF$ a natural transformation such that $\FF$ is Karoubi-localising. Then in the commutative square
\[
\begin{tikzcd}
\F \ar[r]\ar[d] & \FF \ar[d, "\simeq"] \\
\F^{\cotower{-n}} \ar[r]\ar[ur,dotted] & \FF^{\cotower{-n}}
\end{tikzcd}
\]
the right vertical map is an equivalence by Proposition~\reffour{proposition:T-of-localising}, and hence the dotted lift exists and is essentially unique. In particular, any map $\F \Rightarrow \FF$ with a Karoubi-localising target yields an augmentation 
\[
\F \Rightarrow \F^{\cotower{0}} \Rightarrow \F^{\cotower{-1}} \Rightarrow \cdots \Rightarrow \F^{\cotower{-n}} \Rightarrow \cdots \Rightarrow \FF
\]
of the decoration sequence of $\F$, and so induces a map
\[
\F^{\cotower{-\infty}} \Rightarrow \FF .
\]

\begin{proposition}
\label{proposition:connective-covers}%
Let $\FF\colon \Catp \to \Spa$ be a Karoubi-localising functor and let $\F \Rightarrow \FF$ be a map which exhibits $\F$ as the pointwise $m$-connective cover of $\FF$. Then the induced map $\T\F \Rightarrow \T\FF$ exhibits $\T\F$ as the pointwise $(m-1)$-connective cover of $\T\FF \simeq \FF$. 
As a consequence, for $\F = \tau_{\geq m}\FF$ we have $\F^{\cotower{-n}} = \tau_{\geq m-n-1}\F$ and the associated map
$\F^{\cotower{-\infty}} \Rightarrow \FF$ is an equivalence.
\end{proposition}
\begin{proof}
For $(\C,\QF) \in \Catp$ consider the commutative diagram
\[
\begin{tikzcd}
\FF\Hyp(\C)\ar[r,equal]\ar[d]  & \FF\Hyp(\C) \ar[r]\ar[d] & \FF(\Hyp(\C_{\B\NN}))\ar[d] \\
\FF(\C,\QF) \ar[r] & {\B\ZZ} \otimes \FF(\C,\QF) \ar[r, "\ass"] & \FF((\C,\QF)_{\B\ZZ}) \ .
\end{tikzcd}
\]
Since $\FF$ is Karoubi-localising, Theorem~\reffour{theorem:circle-decomposition} implies that the right square is exact. Furthermore, the horizontal arrows in that square admit retractions, and are hence injective on homotopy groups. We conclude that the right square remains exact after applying $\tau_{\geq m}$. The total fibre of the external rectangle then coincides with that of the left square, so the map $\T\F = \T(\tau_{\geq m}\FF) \Rightarrow \T\FF$ can be identified with the map induced on horizontal fibres in the square
\[
\begin{tikzcd}[column sep=5pt]
\tau_{\geq m}\FF(\C,\QF) \ar[rr]\ar[d] && \tau_{\geq m}[{\B\ZZ} \otimes \FF(\C,\QF)] \ar[r,phantom, "="]\ar[d] & \tau_{\geq m}\FF(\C,\QF) \oplus \tau_{\geq m}(\Sig\FF(\C,\QF))\ar[d] \\
\FF(\C,\QF) \ar[rr] && \B\ZZ \otimes \FF(\C,\QF) \ar[r,phantom, "="] & \FF(\C,\QF) \oplus \Sig\FF(\C,\QF) \\
\end{tikzcd}
\]
which is the map $\Om\tau_{\geq m}(\Sig\FF) \Rightarrow \Om\Sig\FF=\FF$. This map is indeed an $(m-1)$-connective cover, as desired.
\end{proof}


\begin{corollary}
\label{corollary:decorated-K}%
For $n \geq 0$ the map $\K^{\cotower{-n}} \Rightarrow \KK$ exhibits $\K^{\cotower{-n}}$ as the $(-n)$-connective cover $\tau_{\geq -n}\KK$ of $\KK$. 
In particular, the induced map 
$\K^{\cotower{-\infty}} \Rightarrow \KK$ is an equivalence.
\end{corollary}
\begin{proof}
We have $\K^{\cotower{0}} = \K((-)^{\natural}) = \tau_{\geq 0}\KK$ by Proposition~\reffour{proposition:T-of-localising} and the definition of $\KK$. The desired result hence follows from Proposition~\reffour{proposition:connective-covers}.
\end{proof}

Let us now bootstrap the above result from $\K$- to Grothendieck-Witt-spectra. 

\begin{proposition}
\label{proposition:decorated-GW}%
For all $n\geq 0$, we have that the square 
\[
\begin{tikzcd} \GW^{\cotower{-n}} \ar[r] \ar[d] & \KGW \ar[d] \\
(\tau_{\geq -n} \KK)^\hC \ar[r] & \KK^\hC
\end{tikzcd}
\]
is exact. In particular, the map $\GW^{\cotower{-\infty}} \Rightarrow \KGW$ is an equivalence.
\end{proposition}
\begin{proof}
Consider the square
\[
\begin{tikzcd}
\GW \ar[r]\ar[d] & \KGW\ar[d] \\
\K^{\hC} \ar[r] & \KK^{\hC} 
\end{tikzcd} \ ,
\]
which is exact by Theorem~\reffour{theorem:Karoubi-pullback-squares}. Applying the functor $(-)^{\cotower{-n}}$ using Proposition~\reffour{proposition:T-of-localising} and Lemma~\reffour{lemma:preserves-limits}\;\reffouritem{item:T-preserves-limits-and-colimits} we obtain the exact square
\[
\begin{tikzcd}
\GW^{\cotower{-n}} \ar[r]\ar[d] & \KGW\ar[d] \\
(\K^{\hC})^{\cotower{-n}} \ar[r] & \KK^{\hC} \ .
\end{tikzcd} 
\]
Corollary~\reffour{corollary:decorated-K} together with the fact that the operation $\T(-)$ preserves all limits by Lemma~\reffour{lemma:preserves-limits}~\reffouritem{item:T-preserves-limits-and-colimits} identifies this with the square from the statement. Finally, taking the colimit over $n \in \mathbb N$, the bottom horizontal map becomes the map
\[
\colim_{n \geq 0}(\KK^{\hC})^{\cotower{-n}} = \colim_{n \geq 0}(\tau_{\geq -n} \KK)^{\hC} \longrightarrow \KK^{\hC},
\]
which is an equivalence by \cite{NS}*{Lemma I.2.6 (ii)} and hence the top horizontal arrow is an equivalence as well.
\end{proof}



\begin{proposition}
\label{proposition:decorated-L}%
For all $n\geq 0$, we have that the square 
\[
\begin{tikzcd} \L^{\cotower{-n}} \ar[r] \ar[d] & \KL \ar[d] \\
(\tau_{\geq -n} \KK)^\tC \ar[r] & \KK^\tC
\end{tikzcd}
\]
is exact. In particular, the induced map $\L^{\cotower{-\infty}} \Rightarrow \LL$ is an equivalence.
\end{proposition}
\begin{proof}
The proof proceeds as that of Proposition~\reffour{proposition:decorated-GW} by applying the functor $(-)^{\cotower{-n}}$, this time to the exact square
\[
\begin{tikzcd}
\L \ar[r]\ar[d] & \KL\ar[d] \\
\K^{\tC} \ar[r] & \KK^{\tC} 
\end{tikzcd} \ ,
\]
equally established in Theorem~\reffour{theorem:Karoubi-pullback-squares}. The resulting square can then identified with the square in the proposition by using Proposition~\reffour{proposition:T-of-localising}, Proposition~\reffour{corollary:decorated-K} and the fact that 
$(-)^{\cotower{-n}}$ 
commutes with the Tate construction by 
Lemma~\reffour{lemma:preserves-limits}\;\reffouritem{item:T-preserves-limits-and-colimits}. The last part follows again from \cite{NS}*{Lemma I.2.6 (ii)}.
\end{proof}


\begin{remark}
\label{remark:all-verdier}%
Corollary~\reffour{corollary:decorated-K} and Propositions \reffour{proposition:decorated-GW} and \reffour{proposition:decorated-L} imply that the natural transformation $\F \Rightarrow \F^{\cotower{-\infty}}$ exhibits $\F^{\cotower{-\infty}}$ as the initial Karoubi-localising functor under $\F$ for $\F=\K,\GW$ and $\L$, respectively. In Theorem~\reffour{theorem:karoubification}
below 
we will show that this in fact holds for any Verdier-localising functor $\F\colon \Catp \to \Spa$.
\end{remark}




\begin{remark}
\label{rem:spectral_rothenberg_lower_decoration}%
Propositions \reffour{proposition:decorated-GW} and \reffour{proposition:decorated-L}
give rise to exact squares
\[
\begin{tikzcd}
	\GW^{\cotower{-n+1}} \ar[r] \ar[d] & \GW^{\cotower{-n}} \ar[d] && \L^{\cotower{-n+1}} \ar[r] \ar[d] & \L^{\cotower{-n}} \ar[d] \\
	(\tau_{\geq -n+1}\KK)^{\hC} \ar[r] & (\tau_{\geq -n}\KK)^{\hC} && (\tau_{\geq -n+1}\KK)^{\tC} \ar[r] & (\tau_{\geq -n}\KK)^{\tC}
\end{tikzcd}
\]
and consequently, one obtains fibre sequences
\[
\begin{tikzcd}[row sep=tiny]
	\GW^{\cotower{-n+1}}(-) \ar[r] & \GW^{\cotower{-n}}(-) \ar[r] & \mathbb{S}^{-n} \otimes \KK_{-n}(-)^{\hC}\,, \\
	\L^{\cotower{-n+1}}(-) \ar[r] & \L^{\cotower{-n}}(-) \ar[r] & \mathbb{S}^{-n} \otimes  \KK_{-n}(-)^{\tC}
\end{tikzcd}
\]
the second of which on homotopy groups induces Ranicki's Rothenberg sequences for lower L-theory \cite{Ranicki-lower}*{Theorem 17.2}. 
\end{remark}






\subsection{Karoubi approximation via Day cotensoring}
\label{subsection:karoubi-approximation-verdier}%



In this subsection we turn to the behaviour of the ultimate decoration construction for general Verdier-localising functors, and aim to show that it always provides an initial Karoubi-localising approximation. The same in fact goes for the Lattice-Tate construction used in the proof of Proposition~\reffour{proposition:localising-are-stable}, and our approach is to consider both of them as instances of Day cotensoring (the right adjoint to Day convolution), and then formulate sufficient conditions under which iterative Day cotensoring leads to Karoubi approximation. Let us begin by setting up the stage.

In what follows, recall that a functor $\C \to \D$ among locally small $\infty$-categories is said to be small if it is left Kan extended from a small full subcategory $\C_0 \subseteq \C$. Given a locally small $\infty$-category $\C$, its Ind-completion $\Ind(\C)$ is by definition the smallest full subcategory of $\Fun(\C\op,\Sps)$ containing the representable presheaves and closed under (small) filtered colimits. In particular, $\Ind(\C)$ is again locally small. If $\C$ admits finite colimits then we can also identify $\Ind(\C)$ with the $\infty$-category of finite-limit-preserving small presheaves $\C\op \to \Sps$. 


\begin{construction}
\label{const:day-conv}%
Let $\cB$ be a locally small and idempotent complete symmetric monoidal $\infty$-category acting on an idempotent complete locally small $\infty$-category $\E$, and assume that $\cB$ and $\E$ admit finite colimits compatibly with the tensor product and action. The $\infty$-categories $\Ind(\cB)$ and $\Ind(\E)$ are then cocomplete and locally small, and $\Ind(\cB)$ inherits a symmetric monoidal structure and an action on $\Ind(\E)$, both compatible with small colimits. 

By Corollary~\reffour{corollary:persist-locally-small}, the $\infty$-category $\Fun^{\sm}(\Catp,\Ind(\cB))$ of small functors $\Catp \to \Ind(\cB)$ admits a symmetric monoidal structure by Day convolution, and acts via Day convolution on $\Fun(\Catp,\Ind(\E))$, where in the latter we do not have to restrict attention to small functors. Here we recall that the Day convolution of a small functor $\T\colon \Catp \to \Ind(\cB)$ and a functor $\F$ from $\Catp$ to either $\Ind(\cB)$ or $\Ind(\E)$ is given by the left Kan extension along $\otimes\colon \Catp \times \Catp \to \Catp$ of the functor $((\C,\QF),(\D,\QFD)) \mapsto \T(\C,\QF)\otimes \F(\C,\QF)$, and the condition that $\T$ is small and $\Catp$ is closed ensures that this left Kan extension is well-defined.

Let us write $\Fun^{\omega}(\Catp,\Ind(\cB)) \subseteq \Fun^{\sm}(\Catp,\Ind(\cB))$ for the full subcategory spanned by the compact objects in $\Fun^{\sm}(\Catp,\Ind(\cB))$. We claim that a small $\T\colon \Catp \to \Ind(\cB)$ is compact if and only if it is a retract of a functor of the form
\[
\colim_{i \in \I}\Gamma(i)_!\tau(i) 
\]
for some finite $\infty$-category $\I$, some functor $\Gamma\colon \I\op \to \Catp$, and some functor $\tau\colon \I \to \cB$, where $\Gamma(i)_!$ means left Kan extension along the functor $\ast \to \Catp$ picking out $\Gamma(i)$. 
Indeed, on the one hand, any such functor is compact, being a retract of a finite colimit of compacts. On the other hand, if $\T$ is a compact small functor then pick a %
small full subcategory $\mathcal{U} \subseteq \Catp$ from which $\T$ left Kan extends, so that we can write $\T$ as a small colimit of the form
\[
\T = \colim_{\stackrel{[(\C,\QF) \to (\D,\QFD)]}{\in \Twar(\mathcal{U})}}(\D,\QFD)_!\T(\C,\QF) .
\]
By~\cite{HTT}*{Variant 4.2.3.15} there exists a poset $\J$ and a cofinal map $\J \to \Twar(\mathcal{U})$, so that the above colimit can be computed after restricting to $\J$. %
Writing $\J$ as a colimit of finite subposets $\I \subseteq \J$ and using the compactness of $\T$ we obtain that
$\T$ is a retract of a functor of the form 
\[
\colim_{i \in \I}\Gamma(i)_!\bar{\tau}(i)
\]
for some finite poset $\I$, some functor $\Gamma\colon \I\op \to \Catp$, and some functor $\bar{\tau}\colon \I \to \Ind(\cB)$. Finally, note that the above expression is colimit preserving in $\bar{\tau}$; writing the latter as a filtered colimit of diagrams which take values in $\cB$ using~\cite{HTT}*{Proposition 5.3.5.15} %
and invoking again the compactness of $\T$ we may assume without loss of generality that $\bar{\tau}$ takes values in $\cB$, and so we get the claim.

Now since $(\Spaf,\QF^{\uni})_!\one_{\cB}$ is compact and $(\C,\QF)_!b \otimes_{\Day} (\C',\QF')_!b' = ((\C,\QF)\otimes(\C',\QF'))_!(b\otimes b')$ we conclude that the collection of compact functors contains the Day unit and is closed under Day convolution, therefore inheriting from $\Fun^{\sm}(\Catp,\Ind(\cB))$ a symmetric monoidal structure and an action on $\Fun(\Catp,\Ind(\E))$. In addition, since $(\C,\QF)_!b \otimes_{\Day} \F = b\otimes \F(\Funx((\C,\QF),-))$ and $\E$ is idempotent complete we see that the Day action of compact functors preserves $\Fun(\Catp,\E) \subseteq \Fun(\Catp,\Ind(\E))$. All in all, this construction results in a symmetric monoidal $\infty$-category $\Fun^{\omega}(\Catp,\Ind(\cB))$ acting on the $\infty$-category $\Fun(\Catp,\E)$.
\end{construction}


In what follows, we would like to construct Karoubi approximations to Verdier-localising functors $\Catp \to \E$, where $\E$ is as in Construction~\reffour{const:day-conv} for some $\cB$. We will hence fix the assumption that $\E$ has finite limits. We recall that Verdier-localising functors are by definition reduced, that is, send the zero Poincaré $\infty$-category to the terminal object $\ast \in \E$, and hence uniquely lift to functors $\Catp \to \E_{\ast/}$. Replacing $\E$ with $\E_{\ast/}$, we may as well assume that $\E$ is pointed, in which case we may as well assume that $\cB$ is pointed as well. 
We then write $\Fun_{\ast}(\Catp,\E) \subseteq \Fun(\Catp,\E)$ for the full subcategory spanned by the reduced functors.
This full subcategory is reflective, with left adjoint given by $\F^{\red} = \cof[(0,0)_!\F(0) \to \F(-)]$. Similarly, let $\Fun_{\ast}^{\omega}(\Catp,\Ind(\cB)) \subseteq \Fun^{\omega}(\Catp,\Ind(\cB))$ be the full subcategory spanned by the compact reduced functors. 
Given a functor of the form $(\C,\QF)_!b$ as above, its reduction is a finite colimit of such, and in particular compact. We conclude that the reduction preserves compact functors, and so $\Fun_{\ast}^{\omega}(\Catp,\Ind(\cB))$ is again a reflective subcategory of $\Fun^{\omega}(\Catp,\Ind(\cB))$.
Using that $(0,0) \otimes (\D,\QFD) = (0,0)$ for every $(\D,\QF)$ we then get that this reflective localisation is multiplicative with respect to Day convolution (see Proposition~\reffour{proposition:multiplicative-day}), and so we have an induced symmetric monoidal structure on $\Fun_{\ast}^{\omega}(\Catp,\Ind(\cB))$ and an induced action of $\Fun_{\ast}^{\omega}(\Catp,\Ind(\cB))$ on $\Fun_{\ast}^{\omega}(\Catp,\E)$. 
The unit $\one^{\red}$ of $\Fun_{\ast}^{\omega}(\Catp,\Ind(\cB))$ is given by $((\Spaf,\QF^{\uni})_!\one_{\cB})^{\red} = \Poinc \tilde{\otimes} \one_{\cB}$, where $\Poinc$ is viewed as a functor to pointed spaces and $\tilde{\otimes}$ stands for the action of pointed spaces on $\Ind(\cB)$ afforded by the fact that it is pointed with finite colimits. More generally, $((\C,\QF)_!b)^{\red} = \Map((\C,\QF),-)\tilde{\otimes} b$, where $\Map((\C,\QF),-)$ is viewed as pointed by the zero Poincaré functor.

Given the above, we adopt the following assumption:

\begin{assumption}
\label{assumption:cotensor}%
Let $\cB$ and $\E$ be as in Construction~\reffour{const:day-conv} and assume in addition
\begin{enumerate}
\item
$\E$ and $\cB$ are pointed.
\item
$\E$ admits finite limits.
\item
For every $b \in \cB$ the action $b\otimes(-)\colon \E \to \E$ admits a right adjoint $(-)^b\colon \E \to \E$.
\end{enumerate}
\end{assumption}

Assumption~\reffour{assumption:cotensor} ensure that the Day convolution action of $\Fun^{\omega}_{\ast}(\Catp,\Ind(\cB))$ on $\Fun_{\ast}(\Catp,\E)$ is closed, that is, for any reduced compact functor $\T\colon \Catp \to \Ind(\cB)$ the reduced Day convolution action $\T\otimes_{\Day}(-)\colon \Fun_{\ast}(\Catp,\E) \to \Fun_{\ast}(\Catp,\E)$ admits a right adjoint. 
For example if $\T = \colim_{i \in \I}\Gamma(i)_!\tau(i)$ for some finite $\infty$-category $\I$ and functors $\Gamma\colon \I\op \to \Catp$ and $\tau\colon \I \to \cB$ then 
\[
\G^{\T}(\C,\QF) = \lim_{i \in \I\op}\G(\Gamma(i))^{\bet(i)} \in \E,
\]
and for a general $\T$ we will get a retract of such a formula.

\begin{definition}
Let $\T\colon \Catp \to \Ind(\cB)$ be a compact reduced functor equipped with a reduced augmentation, that is, a map $\eps\colon \T \Rightarrow \one^{\red} =  \Poinc\tilde{\otimes} \one_{\cB}$ to the reduced unit of Day convolution. We will say that $(\T,\eps)$ is \defi{Karoubi preserving} if for any Karoubi-localising functor $\G\colon \Catp \to \E$ the induced map $\eps^*\colon \G \Rightarrow \G^{\T}$ is an equivalence. 
\end{definition}


\begin{example}
\label{example:lattice-tate}%
Take $\cB = \Sps^{\fin}_*$ be the $\infty$-category of finite pointed spaces and $\E$ any pointed idempotent complete $\infty$-category with finite limits and colimits, so that $\E$ is canonically a $\cB$-module and the action of $\cB$ on $\E$ is closed.
Consider the reduced compact functor
\[
\T_{\Tate} := \cof[\Map(\Tate^{\omega_1}(\Spaf, (\QF^{\uni})\qshift{1}),-) \to \Map(\Latt^{\omega_1}(\Spaf, \QF^{\uni}),-)],
\]
where the mapping spaces in $\Catp$ are viewed as pointed via the zero Poincaré functor.
Reduced augmentations $\eps\colon\T_{\Tate} \Rightarrow \Poinc$ then correspond to Poincaré functors 
\[
(\Spaf,\QF^{\uni}) \to \fib\bigl[\Latt^{\omega_1}(\Spaf, \QF^{\uni}) \to \Tate^{\omega_1}(\Spaf, (\QF^{\uni})\qshift{1})\bigr] ,
\]
and we pick the one determined by the Lattice-Tate sequence (as in the proof of Proposition~\reffour{proposition:localising-are-stable}):
\[
(\Spaf,\QF^{\uni}) \to \Latt^{\omega_1}(\Spaf, \QF^{\uni}) \to \Tate^{\omega_1}(\Spaf, (\QF^{\uni})\qshift{1}).
\]
Then for any reduced $\F\colon \Catp \to \E$ we have 
\[
\G^{\T_{\Tate}}(\C,\QF) = \fib\big[\G((\C,\QF)\otimes \Latt^{\omega_1}(\Spaf, \QF^{\uni})) \to \G((\C,\QF) \otimes \Tate^{\omega_1}(\Spaf, (\QF^{\uni})\qshift{1}))\big],
\]
and if in addition $\E$ is additive and $\G$ preserves finite direct sums this simplifies to $\G^{\T_{\Tate}}(\C,\QF) = \Om\G((\C,\QF) \otimes \Tate^{\omega_1}(\Spaf, (\QF^{\uni})\qshift{1}))$. 
Since the above Poincaré-Verdier sequence remains at least a Karoubi-Verdier sequence after tensoring with any Poincaré $\infty$-category $(\C,\QF)$ we conclude that $\T_{\Tate}$ is Karoubi preserving.
\end{example}


\begin{example}
\label{example:decoration}%
Take $\cB = \Spa^{\fin}$ be the $\infty$-category of finite spectra and $\E$ any idempotent complete stable $\infty$-category, so that $\E$ is canonically a $\cB$-module and the action of $\cB$ on $\E$ is closed.

Let $\T_{\dec}\colon \Catp \to \Spa$ be given by the desuspension of the total cofibre of 
\[
\begin{tikzcd}
\Sig^{\infty}\Map((\Spaf,\QF^{\uni})_{\B\ZZ},-) \ar[r]\ar[d] & \Sig^{\infty}\Map(\Hyp(\Spaf_{\B\NN}),-) \ar[d] \\
\Sigma^{\infty}\Map((\Spaf,\QF^{\uni}),-)\ar[r] & \Sigma^{\infty}\Map(\Hyp(\Spaf),-),
\end{tikzcd}
\]
where again all mapping spaces in $\Catp$ are canonically viewed as pointed spaces and $\Sigma^{\infty}$ is the reduced suspension spectrum. Equivalently, $\T_{\dec}$ is the fibre of the induced map on vertical cofibres in the above square.
We then define a reduced augmentation $\eps\colon \T_{\dec} \Rightarrow \Sig^{\infty}\Poinc$ by the composite
\begin{align*}
\T &\Rightarrow \cof\big[\Sig^{\infty}\Map((\Spaf,\QF^{\uni})_{\B\ZZ},-) \Rightarrow \Sig^{\infty}\Map((\Spaf,\QF^{\uni}),-)\big] \\
& \Rightarrow \cof\big[(\Sig^{\infty} \Map_*((\Spaf,\QF^{\uni}),-))^{\B\ZZ} \Rightarrow \Sig^{\infty}\Map((\Spaf,\QF^{\uni}),-)\big] \\
& = \Sig^{\infty}\Map((\Spaf,\QF^{\uni}),-)\\
& = \Sig^{\infty}\Poinc.
\end{align*}
For any reduced $\F\colon \Catp \to \E$ the functor
\[
\G^{\T}(\C,\QF) = \Sig\TotFib\left[
\begin{tikzcd}
\G(\Hyp(\C)) \ar[r]\ar[d] & \G(\Hyp(\C_{\B\NN})) \ar[d] \\
\G(\C,\QF) \ar[r] & \G((\C,\QF)_{\B\ZZ})
\end{tikzcd}\right] = \T\G(\C,\QF)
\]
coincides with the decoration construction of Definition~\reffour{definition:decoration} such that the map $\G \Rightarrow \G^{\T}$ identifies with the map $\G\Rightarrow \T\G$ from loc.\ cit., and is therefore an equivalence whenever $\G$ is Karoubi-localising by Proposition~\reffour{proposition:T-of-localising}. In particular, $\T_{\dec}$ is Karoubi preserving.
\end{example}

\begin{lemma}
\label{lemma:T-Karoubi-examples}%
For $\cB$ and $\E$ be as in Example~\reffour{example:lattice-tate} we have $\G^{\T_{\Tate}} = \G(-^{\natural})$ for every Verdier-localising $\G\colon \Catp \to \E$. Similarly, 
For $\cB$ and $\E$ be as in Example~\reffour{example:decoration} we have $\G^{\T_{\dec}} = \G(-^{\natural})$ for every Verdier-localising $\G\colon \Catp \to \E$.
\end{lemma}
\begin{proof}
The claim for $\T_{\mathrm{dec}}$ holds by Proposition~\reffour{proposition:T-of-localising}.
For $\T_{\Tate}$, it will suffice to show that for any Poincaré $\infty$-category $(\C,\QF)$ the kernel of the Verdier projection 
\[
(\C,\QF) \otimes \Latt^{\omega_1}(\Spaf, \QF^{\uni}) \to (\C,\QF) \otimes \Tate^{\omega_1}(\Spaf, (\QF^{\uni})\qshift{1})
\]
is given by $(\C^{\natural},\QF)$. To that end, consider the induced map of Karoubi sequences
\[
\begin{tikzcd}
(\C,\QF)\ar[r]\ar[d] & (\C,\QF) \otimes \Latt^{\omega_1}(\Spaf, \QF^{\uni}) \ar[r]\ar[d] & (\C,\QF) \otimes \Tate^{\omega_1}(\Spaf, (\QF^{\uni})\qshift{1})\ar[d] \\
(\C^{\natural},\QF)\ar[r] & (\C^{\natural},\QF) \otimes \Latt^{\omega_1}(\Spaf, \QF^{\uni}) \ar[r] & (\C^{\natural},\QF) \otimes \Tate^{\omega_1}(\Spaf, (\QF^{\uni})\qshift{1}) 
\end{tikzcd}
\]
in which the bottom one is a fibre sequence in $\Catp$ since its left most term is idempotent complete. Since tensoring with a fixed Poincaré $\infty$-category preserves Karoubi equivalences (see Proposition~\reftwo{lemma:Morita-exact-criterion}\;\reftwoitem{item:preserves-K-equivalences}) 
the vertical maps are Karoubi equivalences, and in particular fully-faithful. The induced map on horizontal kernels in the right square is then equivalent to the base changed Karoubi equivalence
\[
(\C^{\natural},\QF) \times_{(\C^{\natural},\QF) \otimes \Latt^{\omega_1}(\Spaf, \QF^{\uni})}[(\C,\QF) \otimes \Latt^{\omega_1}(\Spaf, \QF^{\uni})] \hookrightarrow (\C^{\natural},\QF).
\]
What we need to show is that this map is actually an equivalence in $\Catp$. Now recall that $\Latt^{\omega_1}(\Spaf,\QF^{\uni})$ is defined as the pairings $\infty$-category 
\[
\Latt^{\omega_1}(\Spaf,\QF^{\uni}) = \Pairings(\Ind^{\omega_1}(\Spaf),\hat{\QF}^{\uni}) ,
\]
and so by Lemma~\refone{lemma:pairing-projection} we have a natural equivalence
\[
(\C,\QF) \otimes \Latt^{\omega_1}(\Spaf,\QF^{\uni}) \simeq \Pairings((\C,\QF) \otimes (\Ind^{\omega_1}(\Spaf),\hat{\QF}^{\uni})) .
\]
Being in particular a right adjoint, $\Pairings$ preserves fibre products. Consequently, we obtain a cartesian square
\[
\begin{tikzcd}
(\C^{\natural},\QF)\times_{\Pairings((\C^{\natural},\QF) \otimes (\Ind^{\omega_1}(\Spaf),\hat{\QF}^{\uni}))} [\Pairings((\C,\QF) \otimes (\Ind^{\omega_1}(\Spaf),\hat{\QF}^{\uni}))] \ar[d]\ar[r,hookrightarrow] & (\C^{\natural},\QF) \ar[d] \\
\Pairings\big((\C^{\natural},\QF) \times_{(\C^{\natural},\QF) \otimes (\Ind^{\omega_1}(\Spaf),\hat{\QF}^\uni)} [(\C,\QF) \otimes (\Ind^{\omega_1}(\Spaf),\hat{\QF}^\uni)] \big) \ar[r,hookrightarrow] & \Pairings(\C^{\natural},\QF) 
\end{tikzcd}
\]
in which the horizontal maps are Karoubi equivalences. To show that the top horizontal map is an equivalence it will hence suffice to show that the Karoubi equivalence of hermitian $\infty$-categories 
\[
(\C^{\natural},\QF) \times_{(\C^{\natural},\QF) \otimes (\Ind^{\omega_1}(\Spaf),\hat{\QF}^\uni)} [(\C,\QF) \otimes (\Ind^{\omega_1}(\Spaf),\hat{\QF}^\uni)] \hookrightarrow (\C^{\natural},\QF)
\]
is an equivalence. But this is detected on the level of underlying stable $\infty$-categories since in a Karoubi equivalence the hermitian structure on the domain is restricted from that of the target.  Hence, we are reduced to showing that the dense full subcategory inclusion
\[
\C^{\natural}  \times_{\C^{\natural} \otimes \Ind^{\omega_1}(\Spaf)} [\C \otimes \Ind^{\omega_1}(\Spaf)] \hookrightarrow \C^{\natural}
\]
is essentially surjective.

Let $\NN$ be the set of non-negative integers, $\SS[\NN] \in \Ind^{\omega_1}(\Spaf)$ the $\NN$-indexed direct sum of copies of $\SS$, and $T\colon \SS[\NN] \to \SS[\NN]$ the map induced by the map $(-)+1\colon \NN \to \NN$.
Suppose given a retract diagram $d \to c \to d$ in $\C^{\natural}$ with $c \in \C$ and let $\alp\colon c \to d \to c$ be the corresponding idempotent endomorphism of $c$. We need to show that $d \otimes \SS \in \C^{\natural} \otimes \Ind^{\omega_1}(\Spaf)$ is contained in $\C \otimes \Ind^{\omega_1}(\Spaf)$.
We claim that $d \otimes \SS$ is in fact equivalent to the coequaliser 
\[
e := \coeq\big[\begin{tikzcd}c \otimes \SS[\NN] \ar[r,yshift=-2pt,"\id"']\ar[r,yshift=2pt,"\alp\otimes T"] & c \otimes \SS[\NN] \end{tikzcd}\big] \in \C \otimes \Ind^{\omega_1}(\Spaf) .
\]
To see this, let us decompose $c$ as a direct sum $c = d \oplus d'$ for some $d'\in \C^{\natural}$. Then $\alp\colon c \to c$ decomposes as a direct sum of the identity on $d$ and $0$ on $d'$, yielding a decomposition 
\begin{align*}
e &= \coeq\big[\begin{tikzcd}[ampersand replacement=\&]d \otimes \SS[\NN] \ar[r,yshift=-2pt,"\id"']\ar[r,yshift=2pt,"\id\otimes T"] \& d \otimes \SS[\NN] \end{tikzcd}\big] \oplus \coeq\big[\begin{tikzcd}[ampersand replacement=\&]d' \otimes \SS[\NN] \ar[r,yshift=-2pt,"\id"']\ar[r,yshift=2pt,"0\otimes T"] \& d' \otimes \SS[\NN] \end{tikzcd}\big]\\
&= d\otimes \coeq\big[\begin{tikzcd}[ampersand replacement=\&, column sep=18pt]\SS[\NN] \ar[r,yshift=-2pt,"\id"']\ar[r,yshift=2pt,"T"] \& \SS[\NN] \end{tikzcd}\big]  \\
&= d \otimes \SS
\end{align*}
as claimed.
\end{proof}



\begin{construction}
\label{construction:karoubification}%
Let $\cB,\E$ be as in Assumption~\reffour{assumption:cotensor} and assume in addition that $\E$ admits sequential colimits. Let $\T\colon \Catp \to \Ind(\cB)$ be a compact reduced functor equipped with a reduced augmentation $\eps\colon\T \Rightarrow \one^{\red}$. Iterating the natural transformation $\G \Rightarrow \G^{\T}$ we obtain a sequence
\[
\G \Rightarrow \G^{\T} \Rightarrow \G^{\T\otimes \T} \Rightarrow ... \Rightarrow \G^{\T^{\otimes n}} \Rightarrow ...
\]
and we write $\G^{\T^{\infty}} = \colim_n\G^{\T^{\otimes n}}$ for its colimit.
We note that, as usual with these constructions, there are a-priori $n+1$ different maps $\G^{\T^{\otimes n}} \Rightarrow \G^{\T^{\otimes n+1}}$ one could consider here, indexed by $i \in \{0,...,n\}$, where the $i$'th one is induced by identifying the source with $\G^{\T^{\otimes i}\otimes \one^{\red} \otimes \T^{\otimes n-i}}$ and using the natural transformation $\eps$ on the component $\one^{\red}$. 
Nonetheless, the automorphisms of $\G^{\T^{\otimes n+1}}$ induced by the permutation automorphisms of $\T^{\otimes n+1}$ act transitively on the set of choices, and so the colimit $\G^{\T^{\infty}}$ (as an object of $\Fun(\Catp,\E)_{\G/}$) is in fact independent on the specific choices we make.  
\end{construction}

For instance, we have $\G^{\T_{\mathrm{dec}}^{\infty}} = \G^{\cotower{-\infty}}$ as functors under $\G$. Let us also write $\G^{\Tate} := \G^{\T_{\Tate}^{\infty}}$. We now come to the main result of this subsection:

\begin{theorem}
\label{theorem:karoubification}%
Let $\cB$ and $\E$ be as in Assumption~\reffour{assumption:cotensor} and assume in addition that $\E$ is differentiable (that is, admits finite limits and sequential colimits, and these two commute) and that the cotensoring $(-)^b\colon \E \to \E$ preserves sequential colimits for any $b \in \cB$.
Let $\T\colon \Catp \to \Ind(\cB)$ be a compact reduced functor and $\F\colon \Catp \to \E$ a Verdier-localising functor. If $\T$ is Karoubi preserving and $\F^{\T}$ is Karoubi invariant then $\F^{\T^{\infty}}$ is Karoubi-localising and the map 
\[
\F \Rightarrow \F^{\T^{\infty}}
\]
exhibits its target as initial in Karoubi-localising functors under $\F$.
\end{theorem}

Combining the above theorem with Lemma~\reffour{lemma:T-Karoubi-examples} we conclude:

\begin{corollary}
\label{corollary:tate-dec}%
If $\E$ is a pointed differentiable $\infty$-category with finite colimits and $\F\colon \Catp \to \E$ is a Verdier-localising functor then $\F \Rightarrow \F^{\Tate}$ is the initial Karoubi-localising functor under $\F$. If $\E$ is furthermore stable then the same holds for $\F \Rightarrow \F^{\cotower{-\infty}}$.
\end{corollary}

\begin{corollary}
\label{corollary:karoubification}%
If $\E$ is a pointed differentiable $\infty$-category with finite colimits then the inclusion 
\[
\Funkloc(\Catp,\E) \to \Funvloc(\Catp,\E)
\]
admits a left adjoint $(-)^{\mathrm{kar}}$. In addition, if $R\colon \E \to \E'$ is any functor between pointed differentiable $\infty$-categories with finite colimits such that $R$ commutes with finite limits and sequential colimits then the square
\[
\begin{tikzcd}
\Funkloc(\Catp,\E)\ar[r]\ar[d,"R\circ (-)"] & \Funvloc(\Catp,\E) \ar[d,"R\circ (-)"] \\
\Funkloc(\Catp,\E')\ar[r] & \Funvloc(\Catp,\E')
\end{tikzcd}
\]
is horizontally left adjointable, that is, the canonical map $(R\F)^{\mathrm{kar}} \Rightarrow R(\F^{\mathrm{kar}})$ is an equivalence for all Verdier-localising functors $\F \colon \Catp \to \E$.
\end{corollary}
\begin{proof}
We set $\cB = \Sps^{\fin}_*$ and form Karoubi-localising approximations via $(-)^{\Tate}$ (using Corollary~\reffour{corollary:tate-dec}) to obtain the horizontal left adjoints.
To show the adjointability property we need to show that if $\F\colon \Catp \to \E$ is a Verdier-localising functor then the map $R\circ \F \Rightarrow R\circ \F^{\Tate}$ induced by $\F \Rightarrow \F^{\Tate}$ is initial in Karoubi-localising functors $\Catp \to \E'$ under $R\circ \F$. Note that we do not assume $R$ preserves finite colimits. In particular, $R$ is not a $\cB$-linear, but only lax $\cB$-linear, and similarly the post-composition functor $R\circ (-)\colon \Fun_{\ast}(\Catp,\E) \to \Fun_{\ast}(\Catp,\E)$ is lax $\Fun^{\omega}_{\ast}(\Catp,\Sps_*)$-linear with respect to Day convolution. This means that on the level of Day cotensoring we have an oplax comparison map
\[
R \circ \F^{\T_{\Tate}} \Rightarrow (R\circ \F)^{\T_{\Tate}}
\]
for $\F \in \Fun_{\ast}(\Catp,\E)$.
The condition that $R$ commutes with finite limits implies that this comparison map is an equivalence. The additional assumption that $R$ commutes with sequential colimits now implies that the left facing map in the induced zig-zag
\[
R \circ \F^{\Tate} = R \circ \bigl[\colim_n\F^{\T_{\Tate}^{n}}\bigr] \Leftarrow \colim_n \bigl[R \circ \F^{\T_{\Tate}^{n}}\bigr] \stackrel{\sim}{\Rightarrow} \colim_n[R\circ \F]^{\T_{\Tate}^{n}} = (R \circ \F)^{\Tate}
\]
is an equivalence as well. By Corollary~\reffour{corollary:tate-dec} we conclude that the induced map
\[
R \circ \F \Rightarrow R \circ \F^{\Tate}
\]
exhibits its target as the initial Karoubi-localising functor under the source, as desired.
\end{proof}





\begin{remark}
The proof of Theorem~\reffour{theorem:karoubification} given below works verbatim if we replace $\Catp$ by $\Catx$, and is applicable to the non-hermitian analogues of Examples~\reffour{example:lattice-tate} and~\reffour{example:decoration}. In the former case, one can either work with the underlying stable $\infty$-categories $\Latt^{\omega_1}(\Spaf)$ and $\Tate^{\omega_1}(\Spaf)$, or with their one-sided counterparts $\Ind^{\omega_1}(\Spaf)$ and $\mathrm{Calc}^{\omega_1}(\Spaf) = \Ind^{\omega_1}(\Spaf)/\Spaf$. In all cases the assumptions of Theorem~\reffour{theorem:karoubification} hold for any Verdier-localising functor, as follows from the argument in the proof of Lemma~\reffour{lemma:T-Karoubi-examples}.
In particular, we get that every Verdier-localising functor on $\Catx$ with a finitely cocomplete pointed differentiable target admits a universal Karoubi-localising approximation, and that the latter can be obtained by the Lattice-Tate (or Ind-Calc) construction, or, alternatively, by the decoration construction if the target is stable. 
Related results can also be found in Blumberg-Gepner-Tabuada~\cite{BGT}*{\S 9}.
\end{remark}


The proof of Theorem~\reffour{theorem:karoubification} is obtained by combining Lemma~\reffour{lemma:T-Karoubi-examples} with the following two, essentially independent results:

\begin{proposition}
\label{proposition:is-karoubi}%
Let $\cB$ and $\E$ be as in Theorem~\reffour{theorem:karoubification}, $\T\colon \Catp \to \Ind(\cB)$ a compact reduced functor which is Karoubi preserving and $\F\colon \Catp \to \E$ a Verdier-localising functor.
If $\F^{\T}$ is Karoubi invariant then $\F^{\T^{\infty}}$ is Karoubi-localising.
\end{proposition}

\begin{proposition}
\label{proposition:is-universal}%
Let $\cB$ and $\E$ be as in Theorem~\reffour{theorem:karoubification}, $\T\colon \Catp \to \Ind(\cB)$ a compact reduced functor which is Karoubi preserving and $\F\colon \Catp \to \E$ a functor.
If $\F^{\T^{\infty}}$ is Karoubi-localising then the map $\F \Rightarrow \F^{\T^{\infty}}$ exhibits its target as the Karoubi approximation of its source.
\end{proposition}


\begin{proof}[Proof of Proposition~\reffour{proposition:is-universal}]
This is a well-known argument that goes back to the work of Schwede on symmetric spectra. A proof in the generality we need is recorded in~\cite{motives}*{Proposition 7.1.2}.

Let $\mathrm{Inj}^{\otimes} = \mathrm{Env}(\mathbb{E}_0^{\otimes})$ denote the symmetric monoidal enveloping $\infty$-category of the $\infty$-operad $\mathbb{E}_0^{\otimes}$. Concretely, the objects of $\mathrm{Inj}$ are finite sets, the morphisms are injective maps of such, and the symmetric monoidal structure is given by disjoint union. 
The augmented object $\eps \colon \T \Rightarrow \one^{\red}$ then extends in an essentially unique manner to a symmetric monoidal functor 
\[
\mathrm{Inj}^{\otimes} \to (\Fun^{\omega}_{\ast}(\Catp,\Ind(\cB))\op)^{\Day} \quad\quad I \mapsto \T^{\otimes I} ,
\]
where on the right hand side we have the symmetric monoidal structure induced on opposites by (localised) Day convolution on $\Fun^{\omega}_{\ast}(\Catp,\Ind(\cB))$.
The action of $\Fun^{\omega}_{\ast}(\Catp,\Ind(\cB))\op$ on $\Fun_{\ast}(\Catp,\E)$ via Day cotensoring then restricts to give an action $\beta$ of $\mathrm{Inj}$ on $\Fun_{\ast}(\Catp,\E)$, where $I \in \mathrm{Inj}$ acts on $\Fun_{\ast}(\Catp,\E)$ by $\G \mapsto \G^{\T^{\otimes I}}$. 
Our assumption that cotensors of the form $(-)^b$ for $b \in \cB$ and finite limits commute with sequential colimits in $\E$ then implies $\mathrm{Inj}$ acts on $\Fun_{\ast}(\Catp,\E)$ via sequential-colimit-preserving functors.

A $\G \in \Fun(\Catp,\E)$ is called $\beta$-stable if $\G \to \G^{\T}$ is an equivalence. In particular, every Karoubi-localising $\G$ is $\bet$-stable by our assumptions on $\T$.
Let now $\F\colon \Catp \to \E$ be such that $\F^{\T^{\infty}}$ is Karoubi-localising. Then $\F^{\T^{\infty}}$ is $\beta$-stable and hence $\F$ is $\beta$-semistable in the sense of \cite{motives}*{Proposition 7.1.2}. Loc.\ cit.\ then implies that the map $\F \to \F^{\T^{\infty}}$ is the initial map to a $\beta$-stable object, and in particular to a Karoubi-localising functor.
\end{proof}


\begin{proof}[Proof of Proposition~\reffour{proposition:is-karoubi}]
Since $\T$ is compact there exists a finite $\infty$-category $\I$ and diagrams $\tau\colon \I \to \Ind(\cB)$ and $\Gamma\colon \I\op \to \Catp$ such that $\T$ is a retract of
\[
\T_{\fin} := \colim_{i \in \I} \Gamma(i)_!\tau(i) ,
\]
and so $\G^{\T}$ is a retract of 
\[
\G^{\T_{\fin}} =  \lim_{i \in \I} \G(-\otimes \Gamma(i))^{\tau(i)}
\]
for every reduced $\G\colon \Catp \to \E$. In particular, if $\G$ is Karoubi-invariant then $\G^{\T}$ is again Karoubi-invariant, since tensoring with any Poincaré $\infty$-category preserves Karoubi equivalences by Proposition~\reftwo{lemma:Morita-exact-criterion}\;\reftwoitem{item:preserves-K-equivalences}. 

Now in the case at hand we start with a Verdier-localising functor $\F\colon \Catp \to \E$, in particular reduced, and $\F^\T$ is Karoubi-invariant by assumption. We hence conclude that $\F^{\T^{\otimes n}}$ is reduced and Karoubi-invariant for every $n \geq 1$ and so $\F^{\T^{\infty}}$ is reduced and Karoubi-invariant. To show that $\F^{\T^{\infty}}$ is Karoubi-localising it will hence suffice to show that it sends Poincar\'e-Verdier sequences to fibre sequences. To that end, let $\alp\colon \Del^1\times \Del^1 \to \Catp$ be a Poincaré-Verdier square
\[
\begin{tikzcd}
(\C,\QF)\ar[r]\ar[d] & (\D,\QFD) \ar[d] \\
(\cB,\Pi)\ar[r] & (\E,\QFE)
\end{tikzcd}
\]
and for $n \in \NN \cup \{\infty\}$ write
\[
\Xi(\alp,n) = \bigl[\F^{\T^{\otimes n}}(\C,\QF)  \to \F^{\T^{\otimes n}}(\cB,\Pi)\times_{\F^{\T^{\otimes n}}(\E,\QFE)} \F^{\T^{\otimes n}}(\D,\QFD)\bigr] \in \Ar(\E)
\]
for the arrow whose invertiblity is equivalent to $\F^{\T^{\otimes n}}(\alp)$ being cartesian.
Write $\F^{\T^{\infty}}$ as the colimit of the sequence
\[
\F \Rightarrow \F^{\T} \Rightarrow \F^{\T\otimes \T} \Rightarrow ... \Rightarrow \F^{\T^{\otimes n}} \Rightarrow ...
\]
where in the $n$'th step the map $\F^{\T^{\otimes n}} \Rightarrow \F^{\T^{\otimes n+1}}$ is induced by the map $\F \to \F^{\T}$ via cotensoring with $\T^{\otimes n}$.
Since sequential colimits in $\E$ commute with finite limits by the assumption that $\E$ is differentiable we have $\Xi(\alp,\infty) = \colim_n \Xi(\alp,n)$.
Our goal is hence to show that 
$\colim_n \Xi(\alp,n)$ is an invertible arrow in $\E$, and we will do so by showing that $\Xi(\alp,n) \to \Xi(\alp,n+1)$ factors through an invertible arrow for every $n \geq 0$. 

Now since $\F^{\T}$ is Karoubi invariant the map $\F \Rightarrow \F^{\T}$ factors as a composite $\F \Rightarrow \F(-^{\natural}) \Rightarrow \F^{\T}$, and so the map 
$\F^{\T^{\otimes n}} \Rightarrow \F^{\T^{\otimes n+1}}$ factors through 
$[\F(-^{\natural})]^{\T^{\otimes n}}$. At the same time, the induced map 
$\F^{\T^{\otimes n}} \to [\F(-^{\natural})]^{\T^{\otimes n}}$ is a retract of the induced 
map $\F^{\T^{\otimes n}_{\fin}} \to [\F(-^{\natural})]^{\T^{\otimes n}_{\fin}}$.
To finish the proof it will hence suffice to show that the map of squares
\[
\begin{tikzcd}
\lim\limits_{I \in \I^n} \F((\C,\QF)\otimes \Gamma(I))^{\beta(I)} \ar[r]\ar[d] & \lim\limits_{I \in \I^n} \F((\D,\QFD)\otimes \Gamma(I))^{\beta(I)} \ar[d] \\
\lim\limits_{I \in \I^n} \F((\cB,\Pi)\otimes \Gamma(I))^{\beta(I)}\ar[r] & \lim\limits_{I \in \I^n} \F((\E,\QFE)\otimes \Gamma(I))^{\beta(I)}
\end{tikzcd}
\]
\[
\Downarrow
\]
\[
\begin{tikzcd}
\lim\limits_{I \in \I^n} \F([(\C,\QF)\otimes \Gamma(I)]^{\natural})^{\beta(I)} \ar[r]\ar[d] & \lim\limits_{I \in \I^n} \F([(\D,\QFD)\otimes \Gamma(I)]^{\natural})^{\beta(I)} \ar[d] \\
\lim\limits_{I \in \I^n} \F([(\cB,\Pi)\otimes \Gamma(I)]^{\natural})^{\beta(I)}\ar[r] & \lim\limits_{I \in \I^n} \F([(\E,\QFE)\otimes \Gamma(I)]^{\natural})^{\beta(I)}
\end{tikzcd}
\]
factors through a cartesian square. Now by Lemma~\reffour{lemma:tensor-preserves-karoubi-squares}  each of the squares
\[
\begin{tikzcd}
{[(\C,\QF)\otimes \Gamma(I)]^{\natural}} \ar[r]\ar[d] & {[(\D,\QFD)\otimes \Gamma(I)]^{\natural}} \ar[d] \\
{[(\cB,\Pi)\otimes \Gamma(I)]^{\natural}}\ar[r] & {[(\E,\QFE)\otimes \Gamma(I)]^{\natural}}
\end{tikzcd}
\]
is a Poincaré-Karoubi square among idempotent complete Poincaré $\infty$-categories, hence in particular cartesian. Note that the collection of cartesian squares forms a reflective subcategory inside all squares, with reflection given by replacing the top left corner by the fibre product of the remaining terms. We conclude that the above map of squares factors through the square
\[
\begin{tikzcd}
\lim\limits_{I \in \I^n} \F((\C,\QF)(I))^{\beta(I)} \ar[r]\ar[d] & \lim\limits_{I \in \I^n} \F((\D,\QFD)\otimes \Gamma(I))^{\beta(I)} \ar[d] \\
\lim\limits_{I \in \I^n} \F((\cB,\Pi)\otimes \Gamma(I))^{\beta(I)}\ar[r] & \lim\limits_{I \in \I^n} \F((\E,\QFE)\otimes \Gamma(I))^{\beta(I)},
\end{tikzcd}
\]
where
\[
(\C,\QF)(I) := [(\cB,\Pi) \otimes \Gamma(I)] \times_{(\E,\QFE) \otimes \Gamma(I)} [(\D,\QFD) \otimes \Gamma(I)] .
\]
We now observe that the last square is cartesian since each of the squares
\[
\begin{tikzcd}
(\C,\QF)(I) \ar[r]\ar[d] & (\D,\QFD)\otimes \Gamma(I) \ar[d] \\
(\cB,\Pi)\otimes \Gamma(I)\ar[r] & (\E,\QFE)\otimes \Gamma(I),
\end{tikzcd}
\]
is a Poincaré-Verdier square by Proposition~\reftwo{lemma:Morita-exact-criterion}\;\reftwoitem{item:preserves-PV-proj},
and $\F$ is assumed Verdier-localising.
\end{proof}

It is natural to ask if Theorem~\reffour{theorem:karoubification} can be extended beyond the realm of Verdier-localising functors.
To this end, for $\cB,\E$ and $\T$ as in Theorem~\reffour{theorem:karoubification}, let us denote by $\cK_\T$ the class of functors $\F\colon \Catp \to \E$ for which the map $\F \Rightarrow \F^{\T^{\infty}}$ is a Karoubi-localising approximation. 
By Proposition~\reffour{proposition:is-universal}, this is exactly the class of functors for which $\F^{\T^{\infty}}$ is Karoubi-localising, which means that the class $\cK_\T$ is closed under all operations which commute with $(-)^{\T^{\infty}}$ and preserves Karoubi-localising functors, notably: 
\begin{enumerate}
\item
Pre-composition with $(\C,\QF)\otimes (-)$ for some $(\C,\QF) \in \Catp$.
\item
Finite limits.
\item
All colimits that commute with finite limits in $\E$. 
\end{enumerate}
Now in the case of $\T=\T_{\Tate},\T_{\dec}$, the class $\cK_\T$ contains all Verdier-localising functors by Theorem~\reffour{theorem:karoubification}
and Lemma~\reffour{lemma:T-Karoubi-examples}, and hence also contains all functors of the form $\F((\C,\QF)\otimes (-))$ for Verdier-localising $\F$. We note that $\F((\C,\QF)\otimes (-))$ are not Verdier-localising in general (see Remark~\reftwo{remark:no-verdier}). 

\begin{corollary}
\label{corollary:translated}%
For any Poincaré $\infty$-category $(\C,\QF)$ the maps
\[
\GW((\C,\QF)\otimes(-)) \to \KGW((\C,\QF)\otimes (-)) \quad\text{and}\quad \L((\C,\QF)\otimes (-)) \to \KL((\C,\QF)\otimes(-))
\]
exhibit their targets as initial Karoubi-localising functors under the sources. 
\end{corollary}

We note that the statement of Corollary~\reffour{corollary:translated} is not a formal consequence of the analogous statement for $\KGW$ and $\KL$:
While pre-composition with $(\C,\QF)\otimes (-)$ preserves Karoubi-localising functors, there is a-priori no reason for it to preserve Karoubi-localising approximations, and we do not know if this is true in general. 

Corollary~\reffour{corollary:translated} provides a non-trivial class of non-Verdier-localising additive functors for which a Karoubi-localising approximation can be obtained via the Tate or decoration construction. This does not hold for all additive functors, as the following example shows:

\begin{example}
\label{example:not-all-additive}%
Let $\cB,\E$ be as in Assumption~\reffour{assumption:cotensor} with $\E = \Spa$, and let $\T\colon \Catp \to \Ind(\cB)$ be a Karoubi preserving compact reduced functor (e.g., we can take $\cB$ and $\T$ to be as in Example~\reffour{example:lattice-tate}, or as in Example~\reffour{example:decoration}). In particular, $\T$ is a retract of $\colim_{i \in \I}\Gamma(i)_!\tau(i)$ for some finite $\infty$-category $\I$ and functors $\tau\colon \I \to \cB$ and $\Gamma\colon \I\op \to \Catp$.
Then we can find a large enough uncountable cardinal $\kappa$ such that $\Gamma$ takes values in Poincaré $\infty$-categories with $\kappa$-small mapping spectra.
Let $R = \QQ[x_i]_{i \in \I}$ be the polynomial ring over the rational numbers with set of indeterminates $\{x_i\}_{i \in I}$ of cardinality $|I| = \kappa$, $K$ be the fraction field of $R$, and set 
\[
\F := \KGW(\Funx((\Dperf(K),\QF^{\sym}),-))\colon \Catp \to \Spa,
\]
so that $\F$ is an additive functor. 
We claim that $\F^{\T^{\infty}}$ is not Karoubi-localising. More precisely, we claim that $\F^{\T^{\infty}}$ sends the Karoubi sequence
\[
(\Dperf_K(R),\QF^{\sym}) \to (\Dperf(R),\QF^{\sym}) \to (\Dperf(K),\QF^{\sym})
\]
of Corollary~\reftwo{corollary:ore-localisation} to a sequence of the form $0 \to 0 \to E$ for some non-zero spectrum $E$. To see this, we first note that if $\A$ is a stable $\infty$-category whose mapping spectra are all $\kappa$-small then every exact functor $\Dperf(K) \to \A$ is zero. Indeed, if we let $x \in \A$ be the image of $K \in \Dperf(K)$ and $x$ is not zero then the induced map $K \to \pi_0\map_{\A}(x,x)$ is injective and this is not possible since $K$ is not $\kappa$-small. Secondly, if $\A$ is a stable $\infty$-category and $M \in \Dperf(R)\otimes \A$ is an object whose image in $\Dperf(R')\otimes \A$ is zero for every ring homomorphism $R \to R'$ to a countable ring $R'$, then $M=0$. Indeed, $R$ is a filtered colimit of the countable polynomial rings $\QQ[x_{i_1},...,x_{i_n}]$ ranging over the finite subsets $\{i_1,...,i_n\} \subseteq I$, and so $\Dperf(R)$ is a filtered colimit of the various $\Dperf(\QQ[x_{i_1},...,x_{i_n}])$, and similarly $\Dperf(R)\otimes \A$ is a filtered colimit of the various 
$\Dperf(\QQ[x_{i_1},...,x_{i_n}]) \otimes \A$; the claim then follows from the fact that each of the inclusions $\QQ[x_{i_1},...,x_{i_n}] \subseteq R$ admits a retraction (by sending the other variables to zero, for example).
Combining the above two points with the fact that the mapping spectra in $\Dperf(R')\otimes \A$ are $\kappa$-small whenever $R'$ is countable and $\A$ has $\kappa$-small mapping spectra (indeed, it is enought to consider pure tensors, in which case the mapping spectra are tensor products of a $\kappa$-small spectrum and a countable spectrum, hence again $\kappa$-small) we conclude that if $\A$ has $\kappa$-small mapping spectra then any exact functor $\Dperf(K) \to \Dperf(R)\otimes \A$ is zero, and so $\F^{\T^{\infty}}$ vanishes on $(\Dperf(R),\QF^{\sym})$, and similarly on any Poincaré subcategory of it. We now show that $\F^{\T^{\infty}}(\Dperf(K),\QF^{\sym}) \neq 0$. For this, note that $(\Dperf(K),\QF^{\sym})$ is a symmetric monoidal Poincaré $\infty$-category, so that restriction along the unit Poincaré functor $u\colon (\Spaf,\QF^{\uni}) \to (\Dperf(K),\QF^{\sym})$ gives a natural transformation $u^*\colon \F \Rightarrow \KGW$ fitting in a commutative square
\[
\begin{tikzcd}
\F\ar[r,"u^*"]\ar[d] & \KGW \ar[d,"\simeq"] \\
\F^{\T^{\infty}}\ar[r] & \KGW^{\T^{\infty}} 
\end{tikzcd}
\]
where the right vertical map an equivalence (since we assumed that $\T$ is Karoubi preserving).
It will hence suffice to show that the map
\[
\KGW_0(\Funx((\Dperf(K),\QF^{\sym}),(\Dperf(K),\QF^{\sym}))) \to \KGW_0(\Dperf(K),\QF^{\sym}) ,
\]
induced by the component $u^*_{(\Dperf(K),\QF^{\sym})}$, has a non-zero image. Indeed, the action of $(\Dperf(K),\QF^{\sym})$ on itself via the monoidal product determines a Poincaré section 
\[
(\Dperf(K),\QF^{\sym}) \to \Funx((\Dperf(K),\QF^{\sym}),(\Dperf(K),\QF^{\sym}))
\]
of $u^*_{(\Dperf(K),\QF^{\sym})}$, and $\KGW_0(\Dperf(K),\QF^{\sym})$ is a non-zero ring (it supports a ring homomorphism to $\KK_0(K) = \ZZ$, for example). 
\end{example}

\begin{remark}
\label{rem:accessible-functors-have-Karoubi-approximations}%
We do not know whether every functor $\Catp \to \Spa$ admits a Karoubi-localising approximation. However, in \paperfive 
we will show that if $\kappa$ is an uncountable regular cardinal then every Poincaré-Karoubi square is a $\kappa$-filtered colimits of Poincaré-Karoubi squares among $\kappa$-compact Poincaré $\infty$-categories.
Since $\Catp$ is compactly generated (see Proposition~\reffive{proposition:catp-cath-are-compact}) it then follows from the adjoint functor theorem that any $\kappa$-accessible functor $\Catp \to \Spa$ (hence, eventually, any accessible functor by choosing a large enough $\kappa$) admits a Karoubi-localising approximation. 
\end{remark}





\subsection{The two towers}
\label{subsection:the-two-towers}%

In this final subsection we show that the Tate towers of $\K,\GW$ and $\L$, that is, the towers (or sequences) obtained by iteratively cotensoring with $\T_{\Tate}$, in fact coincide with their decoration towers described in \S\reffour{subsection:universal-dec}. We then revisit the resulting tower from a different perspective, showing that it coincides with the decalage of the stabilisation tower of Definition~\reftwo{definition:stab-construction}.

We begin by proving an analogue of Proposition~\reffour{proposition:connective-covers}:

\begin{proposition}
\label{proposition:connective-covers-tate}%
Let $\FF\colon \Catp \to \Spa$ be a Karoubi-localising functor and let $\F \Rightarrow \FF$ be a map which exhibits $\F$ as the pointwise $m$-connective cover of $\FF$. Then the induced map $\F^{\T_{\Tate}}\Rightarrow \FF^{\T_{\Tate}}$ exhibits $\F^{\T_{\Tate}}$ as the pointwise $(m-1)$-connective cover of $\FF^{\T_{\Tate}} \simeq \FF$. 
As a consequence, for $\F = \tau_{\geq m}\FF$ we have $\F^{\T^{\otimes n}_{\Tate}} = \tau_{\geq m-n-1}\F$ and the associated map
$\F^{\Tate} \Rightarrow \FF$ is an equivalence.
\end{proposition}
\begin{proof}
By definition, the map $\F^{\T_{\Tate}}(\C,\QF)\to \FF^{\T_{\Tate}}(\C,\QF)$ is the map induced on horizontal fibres in the square
\[
\begin{tikzcd}
\F((\C,\QF) \otimes \Latt^{\omega_1}(\Spaf, \QF^{\uni}))\ar[r]\ar[d] & \F((\C,\QF) \otimes \Tate^{\omega_1}(\Spaf, (\QF^{\uni})\qshift{1}) \ar[d] \\
\F((\C,\QF) \otimes \Latt^{\omega_1}(\Spaf, \QF^{\uni}))\ar[r] & \FF((\C,\QF) \otimes \Tate^{\omega_1}(\Spaf, (\QF^{\uni})\qshift{1}) .
\end{tikzcd}
\]
Since $\FF$ is Karoubi-localising it is in particular additive and hence $\F$ is additive as well, so that the terms on the left column vanish by the Eilenberg swindle (as in the proof of Proposition~\reffour{proposition:localising-are-stable}). We may hence identify the map in question with the induced map $\Om\F((\C,\QF) \otimes \Tate^{\omega_1}(\Spaf, (\QF^{\uni})\qshift{1}) \to \FF((\C,\QF) \otimes \Tate^{\omega_1}(\Spaf, (\QF^{\uni})\qshift{1})$ which is indeed an $(m-1)$-connective cover, being the loop of an $m$-connective cover. 
\end{proof}

\begin{corollary}
\label{corollary:two-towers}%
The Tate tower $\F^{\T_{\Tate}^{\otimes n}}$ of $\F$ coincides with the decoration tower $\F^{\<n\>}$ for $\F$ any of the functors $\K,\GW,\L,\K^{\hC},\K^{\tC}$.
\end{corollary}
\begin{proof}
By Lemma~\reffour{lemma:T-Karoubi-examples}
we have $\F^{\T_{\Tate}} \simeq \F^{\T_{\dec}} \simeq \F(-^{\natural})$ for $\F$ any of these three. Now $\K(-^{\natural}) = \tau_{\geq 0}\KK$ and so by Theorem~\reffour{theorem:karoubification} and Proposition~\reffour{proposition:connective-covers-tate} the tower
\[
\K^{\T_{\Tate}} \to \K^{\T_{\Tate}^{\otimes 2}} \to ... \to \K^{\T_{\Tate}^{\otimes n}} \to ... \to \K^{\Tate}
\]
is equivalent to the tower
\[
\tau_{\geq 0}\KK \to \tau_{\geq -1}\KK \to ... \to \tau_{\geq -n+1}\KK \to ... \to \KK ,
\]
and hence coincides with the decoration tower of $\K$ by Corollary~\reffour{corollary:decorated-K}. This holds whether we consider $\K$ as a functor $\Catp \to \Sp$ or $\Catp \to \Sp^{\B\Ct}$. 
Now the Day cotensor operation $\F \mapsto \F^{\T_{\Tate}}$ preserves limits (being right adjoint to the Day convolution with $\T_{\Tate}$), but also colimits by stability, and similarly for $\T_{\dec}$ (as we already observed in Lemma~\reffour{lemma:preserves-limits}). We conclude that the Tate towers of $\K^{\hC}$ and $\K^{\tC}$ are given by
\[
(\tau_{\geq 0}\KK)^{\hC} \to (\tau_{\geq -1}\KK)^{\hC} \to ... \to (\tau_{\geq -n+1}\KK)^{\hC} \to ... \to \KK^{\hC}
\]
and
\[
(\tau_{\geq 0}\KK)^{\tC} \to (\tau_{\geq -1}\KK)^{\tC} \to ... \to (\tau_{\geq -n+1}\KK)^{\tC} \to ... \to \KK^{\tC} ,
\]
and similarly for the corresponding decoration towers (as we already used in the proof of Proposition~\reffour{proposition:decorated-GW}). Finally, for $\GW$ and $\L$ the result follows from the fibre squares
\[
\begin{tikzcd}
\GW(\C,\QF)\ar[r]\ar[d] & \KGW(\C,\QF)\ar[d] & \L(\C) \ar[r] \ar[d] & \KL(\C,\QF) \ar[d] \\
\K(\C,\QF)^\hC\ar[r] & \KK(\C,\QF)^\hC & \K(\C,\QF)^\tC \ar[r] & \KK(\C,\QF)^{\tC}.
\end{tikzcd}
\]
of Theorem~\reffour{theorem:Karoubi-pullback-squares}. In particular, the Tate towers of $\GW$ and $\L$ fit in the same form of fibre squares as we recorded from the decoration tower in Theorem~\reffour{proposition:decorated-GW} and Proposition~\reffour{proposition:decorated-L}.
\end{proof}


We finally offer a third perspective on the equivalent towers of Corollary~\reffour{corollary:two-towers}. Namely, recall from Definition~\reftwo{definition:stab-construction} the two-sided sequence
\[
\dots \longrightarrow \SS^{-1} \otimes \KGW(\C,\QF\qshift{1}) \longrightarrow \KGW(\C,\QF) \longrightarrow \SS^1 \otimes \KGW(\C,\QF\qshift{-1}) \longrightarrow \dots
\]
whose maps are the boundary maps of the Bott-Genauer sequences associated to the various $(\C,\QF\qshift{i})$. By Proposition~\reftwo{proposition:Q-bordism}, the colimit of this sequence is canonically identified with $\KGW^\bord(\C,\QF) \simeq \KL(\C,\QF)$, while the limit is the cobordification $\KGW^\mathrm{cbord}(\C,\QF) \simeq \fib[\LL(\C,\QF) \rightarrow \KK(\C,\QF)^\tC]$. We can therefore regard this sequence as a (non-separated) filtration $\mathrm{stab}_\KGW$ of $\KL$, the stabilisation filtration, and consider its décalage $\dec(\mathrm{stab}_\KGW)$; we will recall the definition in detail below. Roughly, the décalage of a filtered spectrum is a new filtered spectrum, whose associated graded arises from that of the original filtration by taking homology with respect to the canonical differential on it (that also gives rise to the first differential of the associated spectral sequence). In particular, it has the same limit and colimit as the original filtration. 

We claim that the filtration formed by $\L^{\cotower {-n}}(\C, \QF)$ is the décalage of the stabilisation filtration for $n\geq 0$. More precisely, we have:

\begin{proposition}
\label{prop:dec-stab}%
There are cartesian squares
\[
\begin{tikzcd} 
\dec_n(\mathrm{stab}_\KGW) \ar[d] \ar[r] & \LL \ar[d] \\
(\tau_{\geq -n} \KK)^\tC \ar[r] & \KK^\tC
\end{tikzcd}
\]
for all $n \in \mathbb Z$, so that in particular we have $\dec_{n}(\mathrm{stab}_\KGW) \simeq \L^{\cotower{-n}}$ for all $n \geq 0$.
\end{proposition}



To give the definition of the décalage filtration, recall from \cite{ariotta}*{Section 5} the Beilinson $t$-structure on the category $\Fun(\mathbb Z,\Sp)$ of filtered spectra: Its coconnective part is spanned by all $F \colon \mathbb \ZZ \rightarrow \Sp$ for which $F_i$ is $i$-coconnective for all $i \in \mathbb Z$, and its connective part by those $F$ for which $\mathrm{gr}_iF$ is $i$-connective for all $i \in \mathbb Z$, where $\mathrm{gr}_iF = \cof(F_{i-1} \rightarrow F_i)$ is the $i$-th piece of the associated graded spectrum. We will write $\tau_{\geq n}^\mathrm B F$ for the $n$-th connective cover of a filtered spectrum $F$ in this $t$-structure, and denote by $|\tau_{\geq n}^\mathrm B F|$ its colimit.
The décalage $\dec(F)$ of a filtered spectrum $F$ is then given by the sequence
\[
\dots \longrightarrow |\tau_{\geq -n+1}^\mathrm B F| \longrightarrow |\tau_{\geq -n}^\mathrm B F| \longrightarrow \dots
\]
It derives its relevance from the fact that by the main results of \cite{antieau-dec}, the spectral sequence associated to $\dec(F)$ agrees with that associated to $F$ up to a page shift by $1$, that is the $r$-th page for $\dec(F)$ agrees with the $(r+1)$-st page for $F$. This fundamental property allows one to analyse the spectral sequence associated to $F$ entirely in terms of properties of the décalage construction. 

For our purposes the fundamental examples are the following, originally due to Hedenlund, Krause and the eighth author; they are exploited in the forthcoming \cite{KrauseHedenlundNikolaus} to analyse the Tate spectral sequence in detail.

\begin{example}\
\begin{enumerate}
\item Let $X \in \Sp^\hC$ be a $\Ct$-spectrum. Then we have the filtered spectrum, often known as the Tate filtration of $X$, 
\[
\dots \longrightarrow (\SS^{n\sigma} \otimes X)^\hC \longrightarrow (\SS^{(n+1)\sigma} \otimes X)^\hC \longrightarrow \dots,
\]
with maps induced by the cofibre sequence $(\Ct)_+ \rightarrow \ast_+ \rightarrow \mathrm S^\sigma$ of pointed $\Ct$-spaces by tensoring with $\SS^{n\sigma}$. Its limit vanishes and its colimit indeed identifies with $X^\tC$ by noting that the analogous sequence for orbits has vanishing colimit and is constant for the Tate construction since $(\SS^\sigma \otimes X)^\tC \simeq X^\tC$ as a consequence of the cofibre sequence above (and the Tate construction vanishing on induced spectra).

By inspection, the $k$-th connective cover and truncation of this filtration in the Beilinson $t$-structure are then given by 
\[
\dots \longrightarrow (\SS^{n\sigma} \otimes \tau_{\geq k} X)^\hC \longrightarrow (\SS^{(n+1)\sigma} \otimes \tau_{\geq k} X)^\hC \longrightarrow \dots,
\]
and
\[
\dots \longrightarrow (\SS^{n\sigma} \otimes \tau_{\leq k} X)^\hC \longrightarrow (\SS^{(n+1)\sigma} \otimes \tau_{\leq k} X)^\hC \longrightarrow \dots,
\]
respectively: The lower sequence is evidently $k$-coconnective and on account of the cofibre sequence of $\Ct$-spaces above, the $i$-th associated graded of the upper sequence is simply $\mathbb S^i \otimes \tau_{\geq k} X$, which evidently makes it $k$-connective. 

Thus we learn that the décalage of the Tate filtration is given by the Tate construction of the Whitehead filtration:
\[
\dots \longrightarrow (\tau_{\geq -n+1} X)^\tC \longrightarrow (\tau_{\geq -n} X)^\tC \longrightarrow \dots
\]
\item There is a similar filtration 
\[
\dots \longrightarrow (\SS^{n\sigma} \otimes X)^\gC \longrightarrow (\SS^{(n+1)\sigma} \otimes X)^\gC \longrightarrow \dots
\]
for any genuine $\Ct$-spectrum $X \in \Sp^\gC$, with colimit the geometric fixed points $X^\geofix$ and limit $\fib(X^\gC \rightarrow X^\hC) \simeq \fib(X^\geofix \rightarrow X^\tC)$; the argument is the same as above. Again by direct inspection, the $k$-th connective cover and truncation are given by 
\[
\dots \longrightarrow (\SS^{n\sigma} \otimes \tau_{\geq k} X)^\hC \times_{(\SS^{n\sigma} \otimes X)^\hC} (\SS^{n\sigma} \otimes X)^\gC \longrightarrow (\SS^{(n+1)\sigma} \otimes \tau_{\geq k} X)^\hC \times_{(\SS^{(n+1)\sigma} \otimes X)^\hC} (\SS^{(n+1)\sigma} \otimes X)^\gC \longrightarrow \dots,
\]
and
\[
\dots \longrightarrow (\SS^{n\sigma} \otimes \tau_{\leq k} X)^\hC \longrightarrow (\SS^{(n+1)\sigma} \otimes \tau_{\leq k} X)^\hC \longrightarrow \dots,
\]
respectively, since the upper sequence has the same associated graded as in the previous example. It thus follows that the décalage of this filtration of the geometric fixed points has $n$-th term the pullback
\[
\begin{tikzcd}
\dec_n((\SS^{\bullet\sigma} \otimes X)^\gC) \ar[r] \ar[d] & X^\geofix \ar[d] \\
(\tau_{\geq -n} X)^\tC \ar[r] & X^\tC.
\end{tikzcd}
\]
\end{enumerate}
\end{example}
 
Now, for the proof of Proposition \reffour{prop:dec-stab} it only remains to observe that by Karoubi periodicity the stabilisation filtration of $\KL(\C,\QF)$ is precisely the filtration of the geometric fixpoints described in the previous example for the genuine $\Ct$-spectrum $\KKR(\C,\QF)$, see the discussion after \reftwo{remark:compare-bordification}. In fact, this works for an arbitrary additive functor $\F \colon \Catp \rightarrow \Sp$: Recall that for any such $\F$ and $(\C,\QF)$ in $\Catp$ we produced a genuine $\Ct$-spectrum $\F^\ghyp(\C,\QF)$, such that the stabilisation filtration $\mathrm{stab}_\mathcal F$ given by
\[
\dots \longrightarrow \SS^{n} \otimes \F(\C,\QF\qshift{-n}) \longrightarrow \SS^{n+1} \otimes \F(\C,\QF\qshift{-n-1}) \longrightarrow \dots
\]
is precisely the filtration
\[
\dots \longrightarrow (\SS^{n\sigma} \otimes \F^\ghyp(\C,\QF))^\gC \longrightarrow (\SS^{(n+1)\sigma} \otimes \F^\ghyp(\C,\QF))^\gC \longrightarrow \dots
\]
of $\F^\ghyp(\C,\QF)^\geofix \simeq \F^\bord(\C,\QF)$ considered in the previous example by \reftwo{theorem:periodicity}. We thus find:

\begin{proposition}
Let $\mathcal F \colon \Catp \rightarrow \Sp$ be additive and $(\C,\QF) \in \Catp$. Then we have exact squares
\[
\begin{tikzcd}
\dec_n(\mathrm{stab}_\mathcal F) \ar[r] \ar[d] & \F^\bord \ar[d] \\
(\tau_{\geq -n} \F^\hyp)^\tC \ar[r] & (\F^\hyp)^\tC
\end{tikzcd}
\]
for all $n \in \mathbb Z$.
\end{proposition}

Inserting $\F = \KGW$ and using Proposition~\reffour{proposition:decorated-L} we thus obtain Proposition \reffour{prop:dec-stab}. %
For $\F = \GW$ the decalaged stabilisation filtration
agrees with that of $\KGW$ %
in negative degrees, %
then terminates at $\L$ in degree $0$, and is constant in positive degrees.
For $\F$ bordism invariant both the stabilisation filtration and its décalage are constant, and finally for $\F = \K$, we find that the stabilisation filtration has vanishing connecting maps and the décalage vanishes.
